\chapter[chap:phase2]{Phase 2: Diff und Ausgabe}

\section{Ziel}

Änderungen zwischen Commits, auch gegenüber dem Working-Copy-Commit, werden lesbar angezeigt.

\section{Umfang}

\startitemize
\item Eigene Implementierung des Myers-Diff-Algorithmus mit Tests gegen bekannte Fälle.
Der Algorithmus sucht eine kürzeste Folge von Einfügungen und Löschungen.
\item Nachbearbeitung der Grenzen nach dem Verfahren von Git: Verschieben von Blöcken
und Bewertung von Einrückungen. Sie ändert nur die Lage der Grenzen, nicht die Zahl der
Änderungen.
Eine Ersetzung, also eine Folge aus Löschungen und Einfügungen, bleibt dabei zusammen,
und eine reine Löschung oder Einfügung rückt nicht an sie heran. Das ist einfacher als bei Git.
Die Einrückung zählt in Bytes, ein Tabulator gilt also als ein Zeichen, bei Git als mehrere.
Der Befehl \type{diff} ruft die Nachbearbeitung auf (\in{Abschnitt}[sec:slide]).
\item Gruppierung in Hunks und Darstellung mit Kontextzeilen (\in{Abschnitt}[sec:hunks]).
\item Farbausgabe und Einbindung eines Pagers (\in{Abschnitt}[sec:pager]).
\item Eine schreibgeschützte Terminal-Oberfläche (TUI), \type{diff --tui}, zum Blättern durch
die Dateien und Hunks (ADR-016, \in{Abschnitt}[sec:tui]). Sie liegt in der Crate
\type{stemma-tui} und hängt vom Kern ab, der Kern nicht von ihr.
\item Der Diff wird in Phase 2 nur für die Anzeige berechnet, also auf Risikostufe 1 nach ADR-004.
Die Verwendung im Merge-Pfad, Risikostufe 2, kam in Phase 3 hinzu (\in{Kapitel}[chap:phase3]).
\item In den Hunks erscheinen verschobene Zeilen als Löschung und Einfügung; der Algorithmus
des Diffs erkennt keine Verschiebungen.
Patience und Histogram verankern auf eindeutigen Zeilen, zeigen verschobene
Blöcke aber ebenfalls nicht und brauchen Myers für die Lücken
(Vergleich der Verfahren in \type{docs/research/moved-lines-heuristics.md}).
Eine Markierung der Verschiebung als reine Anzeige, wie \type{--color-moved} in Git,
ändert die Hunks nicht (\in{Abschnitt}[sec:moved]).
\item Binärdateien erkennt der Diff an einem NUL-Byte in den ersten 8000 Bytes, wie Git. Er zeigt
dann eine Zusammenfassungszeile wie \quotation{Binary files differ} statt des Inhalts
(\in{Abschnitt}[sec:binary]).
\item Eine Textdatei über 10 MiB (10\,485\,760 Bytes) auf einer der beiden Seiten vergleicht der
Diff nicht zeilenweise.
Der Myers-Diff mit linearem Speicher begrenzt den Speicher, aber Zeit und Ausgabe
wachsen mit der Datei. Der Diff zeigt dann eine Zusammenfassungszeile wie bei einer Binärdatei
(\in{Abschnitt}[sec:binary]).
\item Die Seiten des Diffs werden über ihre Kennungen benannt, wie \type{op log} sie
anzeigt. Die Arbeitskopie erscheint als \type{working-copy} mit ihrer Snapshot-Kennung,
ohne Zeitangabe, weil sich die Zeit bei jedem Aufruf ändert (ADR-010).
\item Der Diff lässt sich auf Dateien und Verzeichnisse beschränken, wie \type{git diff -- PFAD}.
Ein Pfad wählt eine Datei oder ein Verzeichnis samt allem darunter; Globs und eine Sprache für
Pfadmuster folgen später. Der Pfad gilt relativ zum Arbeitsverzeichnis und wird wie bei
\type{track} umgeformt; er darf das Repository nicht verlassen. Eine Kennung beginnt immer mit
\type{b3:}, darum lassen sich Kennung und Pfad unterscheiden; ein Pfad, der so beginnt,
steht hinter \type{--}. Auch ein Argument mit \type{B3:} gilt als Kennung und wird als ungültige
Kennung abgelehnt, nicht als fehlender Pfad gemeldet.
Steht hinter \type{--} kein Pfad, gilt die Regel ohne \type{--}.
Ein Pfad, der auf keiner der beiden Seiten vorkommt, ist ein Fehler und kein leerer Diff,
damit ein Tippfehler auffällt. Das gilt auch für eine Datei,
die nicht verfolgt wird oder die \type{.stemmaignore} ausschliesst,
denn sie steht in keinem Commit. Der Snapshot
umfasst weiter alle verfolgten Dateien; nur die Ausgabe ist beschränkt, und
\type{--exit-code} betrachtet nur die gewählten Dateien.
\item Ist eine Seite der Working-Copy-Commit, speichert der Diff zuerst einen Snapshot;
hat sich eine verfolgte Datei geändert, schreibt er dazu einen Eintrag ins Operation Log
(ADR-010). Er braucht darum einen Akteur und kann
dieselben Fehler melden wie \type{status}. Mit zwei Kennungen schreibt er nichts.
\item Der Diff liest beide Seiten bei jedem Aufruf neu und legt nichts unter
\type{.stemma/derived/} ab (ADR-017). Ein Cache folgt erst mit einer Abfrage, die ihn braucht.
\item Keine JSON-Form für \type{diff} in Phase 2 (ADR-018). Der Text des Diffs ist für Menschen
bestimmt und kann sich ändern. Ein Skript wertet den Exit-Code aus: \type{--exit-code} liefert 4,
wenn sich eine Datei unterscheidet. Die Hunk-Struktur im Kern besteht aus reinen Daten, auf die
das Format aus Phase 6 aufsetzen kann (\in{Kapitel}[chap:phase6]).
\stopitemize

\section[sec:myers]{Der Myers-Diff}

Der Diff vergleicht zwei Folgen $a$ und $b$ der Längen $N$ und $M$, meist die Zeilen zweier
Dateien. Er sucht eine kürzeste
Folge von Einfügungen und Löschungen, die die alte Folge in die neue überführt. Die Zahl dieser
Änderungen heisst $D$. Myers hat dafür ein Verfahren über den Editiergraphen angegeben
\cite[myers1986ond]. Der Code steht in \type{diff.rs} als \type{diff}.

\subsection{Definition}

Der Editiergraph hat die Gitterpunkte $(x,y)$ mit $0 \le x \le N$ und $0 \le y \le M$.
Ein Schritt von $(x,y)$ nach $(x+1,y)$ löscht das Element $a_x$ der alten Folge.
Ein Schritt nach $(x,y+1)$ fügt das Element $b_y$ der neuen Folge ein.
Eine Diagonale von $(x,y)$ nach $(x+1,y+1)$ kostet nichts und ist nur erlaubt, wenn $a_x = b_y$
gilt. Ein Skript ist ein Weg von $(0,0)$ nach $(N,M)$. Ein kürzestes Skript hat $D$ Schritte,
die keine Diagonale sind. Die Diagonale $k = x - y$ fasst die Punkte mit derselben Differenz
$x - y$ zusammen. Sie liegen auf einer Linie parallel zur Hauptdiagonale. Eine Löschung erhöht
$k$ um eins, eine Einfügung senkt es um eins.

\in{Abbildung}[fig:edit-graph] zeigt den Graphen für die Folgen \type{ABCABBA} und
\type{CBABAC} aus der Arbeit von Myers. Die grauen Diagonalen verbinden gleiche Elemente. Der
dicke Weg ist ein kürzestes Skript mit $D = 5$: zwei Löschungen, eine Einfügung, eine Löschung
und eine Einfügung.

\startplacefigure[title={Editiergraph für \type{ABCABBA} und \type{CBABAC} mit der Mittelschlange},
    reference=fig:edit-graph,location=here]
\externalfigure[images/edit-graph.pdf][width=\textwidth]
\stopplacefigure

\subsection{Invariante}

Die Suche läuft vorwärts vom Anfang und rückwärts vom Ende zugleich. Nach der Runde $d$ kennt
die Vorwärtssuche für jede Diagonale $k$ mit $|k| \le d$ den weitesten erreichten Punkt $V[k]$:
den grössten $x$-Wert, den ein Skript mit $d$ Schritten auf dieser Diagonale erreicht. Die
Rückwärtssuche hält dasselbe vom Ende her fest. Die Runde $d$ berechnet nur die Diagonalen
mit $k \equiv d \pmod 2$. Die anderen behalten den Wert der Runde $d-1$, und ein Skript mit
$d$ Schritten kann auf ihnen nicht enden. Die Felder sind mit null vorbelegt. Die erste Runde
liest deshalb auf den Diagonalen $\pm 1$ den Wert null, und die Suche setzt im Punkt $(0,0)$ an.

\subsection{Komplexität}

Die Suchen besuchen zusammen höchstens $O(D)$ Runden mit je $O(D)$ Diagonalen. Auf einer
Diagonale steigt $x$ über die Runden nur an und sinkt nie. Das Folgen gleicher Elemente kostet
deshalb auf jeder Diagonale zusammen höchstens $\min(N, M)$ Schritte,
auf allen $O(D)$ Diagonalen also
$O(N\,D)$. Die Runden selbst kosten $O(D^2)$. Wegen $D \le N+M$ ist die Laufzeit damit
$O((N+M)\,D)$ \cite[myers1986ond]. Der Speicher ist $O(N+M)$: Die Suche hält zwei Felder der Länge
$2(N+M)+3$, und jeder Aufruf gibt sie wieder frei, bevor er die Teile rekursiv behandelt.

\subsection{Mittelschlange}

Treffen sich die beiden Suchen auf einer Diagonale, liegt dort eine Schlange,
also eine Folge gleicher Elemente entlang
einer Diagonale, auf einem kürzesten Skript.
Das Verfahren zerlegt die Aufgabe an dieser Schlange in zwei Teile und
behandelt jeden Teil genauso. Jeder Teil hat höchstens $\lceil D/2 \rceil$ Schritte, deshalb
endet die Rekursion nach $O(\log D)$ Ebenen. Gemeinsame Präfixe und Suffixe schneidet es vorher ab.

In \in{Abbildung}[fig:edit-graph] liegt die Mittelschlange von $(3,2)$ bis $(5,4)$ auf der
Diagonale $k = 1$. Die Vorwärtssuche hat nach drei Änderungen $(3,2)$ erreicht und läuft die
Schlange bis $(5,4)$ ab. Die Rückwärtssuche vom Ende her erreicht nach zwei Änderungen $(5,4)$
und läuft sie bis $(3,2)$ ab. Zusammen sind das $D = 5$ Änderungen.

\subsection{Korrektheit}

Die Schlange liegt auf einem kürzesten Skript, und die Summe der Schritte der beiden Teile ist
$D$. Ein kürzeres Skript für einen Teil würde das Ganze verkürzen,
also ist jeder Teil ein kürzester Weg.
Ein gemeinsames Element am Rand liegt ebenfalls auf einem kürzesten Skript, deshalb geht beim
Abschneiden nichts verloren \cite[myers1986ond]. Die Teile liegen in der Reihenfolge der Schlange
aneinander. Die so gebaute Folge deckt jedes Element beider Folgen genau einmal ab und hält die
Reihenfolge.

\section[sec:slide]{Nachbearbeitung der Grenzen}

Ein kürzestes Skript ist nicht eindeutig. Dieselben Änderungen lassen sich oft an mehreren Stellen
unterbringen, weil die erste Zeile einer Änderungsgruppe gleich der unveränderten Zeile hinter ihr
sein kann. Die Nachbearbeitung wählt unter diesen Lagen die, die sich am besten liest
(\type{slide} in \type{diff.rs}). Sie folgt dem Verfahren von Git \cite[git-xdiff]. Sie ändert nur
die Lage der Grenzen. Beide Folgen und die Zahl $D$ der Änderungen
bleiben, wie sie sind.
Der Kern bietet sie als \type{diff_lines} für Text und als \type{diff_byte_lines}
für Zeilen aus Bytes an; der Befehl \type{diff} ruft die zweite auf.
Innerhalb einer Gruppe von Änderungen stehen die Löschungen vor den Einfügungen.

Eine Gruppe von Einfügungen oder Löschungen gleitet nach oben oder unten, solange die Zeile, die
sie verlässt, der Zeile gleicht, die sie aufnimmt. \in{Abbildung}[fig:slide] zeigt vier
eingefügte Zeilen. Ein kürzestes Skript kann sie bei der schliessenden Klammer der Funktion a
beginnen lassen (links). Nach dem Gleiten um zwei Zeilen besteht die Gruppe aus der Funktion b
samt der Leerzeile danach (rechts). Unter den möglichen Lagen
wählt die Einrückungsregel, die ihre Gewichte aus dem Code von Git übernimmt: Sie vergleicht die
Einrückung der Zeilen an der Grenze und beachtet Leerzeilen davor und danach.

\startplacefigure[title={Eine Gruppe vor und nach der Nachbearbeitung der Grenzen},
    reference=fig:slide,location=here]
\externalfigure[images/slide.pdf][width=\textwidth]
\stopplacefigure

\subsection{Definition}

Eine Änderungsgruppe ist eine maximale Folge geänderter Zeilen einer Seite. Eine Zeile der alten
Seite ist geändert, wenn das Skript sie löscht, eine der neuen, wenn es sie einfügt. Eine Gruppe
bestehe aus den Zeilen $s$ bis $e-1$ und habe die Länge $g = e - s$. Sie rückt um einen Schritt
nach unten, wenn die Zeile $s$ gleich der Zeile $e$ ist: Dann ist die Zeile $s$ unverändert,
die Zeile $e$ geändert, und die Gruppe belegt die Zeilen $s+1$ bis $e$. Ein Schritt nach oben ist
erlaubt, wenn die Zeile $s-1$ gleich der Zeile $e-1$ ist. Die Lagen einer Gruppe sind alle, die
durch solche Schritte erreichbar sind, ohne eine andere geänderte Zeile zu überschreiten. Mit $u$
Schritten nach oben und $d$ Schritten nach unten sind das die Anfangspositionen von $s-u$ bis
$s+d$.

Jede Lage hat zwei Schnittstellen, vor der ersten und hinter der letzten Zeile der Gruppe.
Für jede liest das Verfahren die Einrückung der Zeile an der Schnittstelle, die Leerzeilen davor
und danach und die Einrückung der nächsten Zeile mit Inhalt auf beiden Seiten, höchstens 20
Leerzeilen weit. Eine Zeile aus Leerraum gilt als leer. Daraus entstehen zwei Zahlen, eine
Strafe $p$ und die wirksame Einrückung $w$, jeweils als Summe über beide Schnittstellen.
Leerzeilen an der Schnittstelle senken die Strafe, der Dateianfang erhöht sie um 1 und das
Dateiende um 21. Das Dateiende zählt zugleich als Leerzeile; seine Strafe sinkt dadurch netto
um 3, und seine wirksame Einrückung ist $-1$.
Springt die Einrückung an der Schnittstelle zurück, kostet das je nach Fall 17 bis 24 Punkte.
Die Gewichte sind die von Git. Zwei Lagen $A$ und $B$ vergleicht der Ausdruck
$60 \cdot \mathrm{sgn}(w_A - w_B) + (p_A - p_B)$; ist er negativ, ist $A$ besser. Die flachere
Einrückung gibt also einen Vorsprung von 60 Punkten, den eine grosse Differenz der Strafen
aufwiegen kann. Bei Gleichstand gewinnt die spätere Lage, also die weiter unten.

Das Verfahren prüft nicht alle Lagen. Das Fenster reicht von der untersten Lage höchstens $g+1$
Schritte nach oben, wie bei Git, und nie weiter als 100. Eine Lage, in der die Gruppe an eine
andere geänderte Zeile grenzte, überspringt es, damit zwei Gruppen nicht verschmelzen.

\subsection{Invariante}

Die beiden Seiten werden getrennt behandelt, und für jede Seite gilt nach jedem Schritt:
Die Gruppe hat weiter $g$ geänderte Zeilen, und die unveränderten Zeilen der Seite ergeben, der
Reihe nach gelesen, dieselbe Folge von Werten wie vorher. Eine Gruppe, die zu einer Ersetzung
gehört, rückt gar nicht. Eine Ersetzung ist eine Folge von Änderungen, die Löschungen und
Einfügungen enthält. Ihre beiden Seiten liegen schon gegeneinander, und das Schieben einer Seite
würde sie auseinanderreissen. Eine reine Löschung oder Einfügung rückt auch nicht an eine
Ersetzung heran.

Die Prüfung, ob die Zeile $s$ der Zeile $e$ gleicht, hält die Invariante. Ohne sie würde ein
Schritt einen Wert aus der Folge der unveränderten Zeilen durch einen anderen ersetzen.

\subsection{Komplexität}

Seien $N$ und $M$ die Längen der beiden Folgen. Die Schritte nach unten und oben vergleichen
Zeilen. Jede Folge unveränderter Zeilen zwischen zwei Gruppen durchläuft das höchstens zweimal,
einmal von der Gruppe darüber nach unten und einmal von der darunter nach oben. Das sind
$O(N+M)$ Vergleiche. Eine Gruppe der Länge $g$ bewertet höchstens $\min(g+1, 100) + 1$ Lagen, und
jede Bewertung liest höchstens 20 Zeilen auf jeder Seite der beiden Schnittstellen. Die Summe der
Längen aller Gruppen ist höchstens $N+M$, also kostet die Bewertung $O(N+M)$ mit dem Faktor 20
als Konstante. Das Markieren der Ersetzungen und der Aufbau des Skripts aus den Marken sind je
ein Durchlauf. Die Laufzeit ist $O(N+M)$ und der Speicher $O(N+M)$ für eine Marke je Zeile.

\subsection{Korrektheit}

Ein Schritt nach unten macht die Zeile $s$ unverändert und die Zeile $e$ geändert. In der Folge
der unveränderten Zeilen der Seite tritt die Zeile $s$ an die Stelle, an der vorher die Zeile $e$
stand. Beide sind gleich, also bleibt die Folge der Werte dieselbe. Ein Schritt nach oben ist das
Spiegelbild. Eine Kette solcher Schritte erhält die Folge damit ebenfalls.

Der Aufbau des Skripts paart die $i$-te unveränderte Zeile der alten Seite mit der $i$-ten der
neuen. Beide Seiten haben gleich viele unveränderte Zeilen, und ihre Wertefolgen waren vor dem
Schieben gleich, weil das Skript sie paarte. Sie sind es auch danach, denn jede Seite hat ihre
Folge behalten. Die Paare sind also gleich, und das Ergebnis ist ein gültiges Skript. Die Zahl
der Änderungen ist die der geänderten Zeilen beider Seiten, und die hat sich nicht geändert. Das
Skript ist damit weiter kürzest.

Dass die Wahl der Lage mit Git übereinstimmt, beweist dieses Argument nicht. Wo der Myers-Diff
und der von Git verschiedene Skripte wählen, weichen auch die Ergebnisse ab. Ein Test in
\type{tests/diff_against_git.rs} vergleicht auf zufälligen Dateien beide Ergebnisse und findet
keinen Unterschied, wo die einfachen Diffs übereinstimmen.

\startinfobox
{\bf Grenzfälle beim Verschieben.}
\startitemize
\item Eine Ersetzung bleibt stehen. Das ist einfacher als bei Git, das die Gruppen beider Seiten
im Einzelnen aufeinander abstimmt.
\item Die Einrückung zählt in Bytes, ein Tabulator also als eines, bei Git als mehrere.
Rückt eine Datei mit Tabulatoren ein, kann das Ergebnis von dem von Git abweichen.
\item Das Dateiende zählt wie eine Leerzeile und hat die wirksame Einrückung $-1$. Es zieht eine
Gruppe am Ende eher nach unten als nach oben; ob sie rutscht, entscheiden die Einrückungen.
\item Ein Skript ohne Änderung hat keine Gruppe, und das Ergebnis ist das Skript selbst.
\item Zwei Gruppen derselben Seite bleiben getrennt, auch wenn eine von ihnen zur anderen hin
rutschen könnte.
\stopitemize
\stopinfobox

\section[sec:hunks]{Hunks und Kontextzeilen}

Ein Skript wird für die Anzeige in Hunks zerlegt. Ein Hunk fasst benachbarte Änderungen mit den
unveränderten Zeilen davor und danach zusammen. Jede Änderung bekommt so viele Kontextzeilen auf
beiden Seiten, wie die Anzeige verlangt. Die Anzahl $c$ setzt der Aufruf; der Befehl \type{diff}
nimmt drei. Die Zerlegung folgt dem Unified-Format von diff und Git. Eine eigene Quelle nennt das
Buch nicht, denn sie ist ein Durchlauf über das Skript. Der Code steht in \type{diff.rs} als
\type{hunks}.

Zwei benachbarte Änderungen landen im selben Hunk, wenn höchstens $2c$ unveränderte Zeilen
zwischen ihnen liegen. Dann überlappen oder berühren sich ihre Kontextfenster. Solche Nachbarn
bilden Ketten, sodass ein Hunk auch Änderungen umfassen kann, die weiter auseinanderliegen.

\in{Abbildung}[fig:hunks] zeigt ein Skript mit drei Änderungen und $c = 3$. Zwischen den ersten
beiden liegen fünf unveränderte Schritte, und $5 \le 2c = 6$ gilt: Ihre Fenster überlappen, sie
bilden einen Hunk. Zwischen der zweiten und der dritten liegen sieben Schritte. Das sind mehr als
sechs, die Fenster berühren sich nicht, und ein neuer Hunk beginnt.

\startplacefigure[title={Gruppierung eines Skripts in Hunks mit $c = 3$},
    reference=fig:hunks,location=here]
\externalfigure[images/hunks.pdf][width=\textwidth]
\stopplacefigure

\subsection{Definition}

Das Skript habe $m$ Schritte, die mit null nummeriert sind. Ein Schritt ist eine Änderung, wenn er
eine Zeile löscht oder einfügt. Zu jeder Änderung $i$ gehört das Fenster der Schritte von $i-c$
bis $i+c$; am Anfang und am Ende des Skripts wird es abgeschnitten. Fenster, die sich überlappen
oder berühren, verschmilzt das Verfahren. Jedes verschmolzene Fenster ist ein Hunk. Die Fenster
zweier benachbarter Änderungen $i < j$ verschmelzen genau dann, wenn $j - i - 1 \le 2c$ gilt.

Ein Hunk speichert die Schritte des Skripts und die Positionen seiner Zeilen: den Anfang und die
Länge je Seite. Der Anfang auf einer Seite ist die Zahl der Zeilen dieser Seite, die die Schritte
davor verbrauchen, und die Länge die Zahl der Zeilen, die die Schritte des Hunks verbrauchen. Die
Zeilen selbst bleiben in den Folgen.

\subsection{Invariante}

Die Hunks folgen der Reihenfolge des Skripts und überlappen sich nicht. Jede Änderung liegt in
genau einem Hunk, und jeder Schritt eines Hunks liegt höchstens $c$ Schritte von einer Änderung
desselben Hunks entfernt. Zwischen zwei aufeinanderfolgenden Hunks liegt mindestens ein Schritt,
der zu keinem gehört, denn Fenster, die sich berühren, hat das Verfahren verschmolzen. Anfang und
Länge eines Hunks zählen auf jeder Seite die Zeilen, die die Schritte davor und die Schritte des
Hunks verbrauchen.

\subsection{Komplexität}

Sei $m$ die Zahl der Schritte, höchstens $N+M$. Das Verfahren durchläuft das Skript einmal, um die
Fenster zu bilden, ein zweites Mal, um die verbrauchten Zeilen je Position zu zählen, und kopiert
die Schritte der Hunks. Jeder Schritt gehört zu höchstens einem Hunk. Laufzeit und Speicher sind
daher $O(m)$, unabhängig von $c$.

\subsection{Korrektheit}

Jede Änderung liegt in ihrem eigenen Fenster, und Verschmelzen verliert keinen Schritt. Also liegt
jede Änderung in einem Hunk. Ein Schritt in einem verschmolzenen Fenster liegt im Fenster einer
Änderung und ist darum höchstens $c$ Schritte von ihr entfernt. Zwei Fenster verschmilzt das
Verfahren, sobald der Anfang des späteren nicht hinter dem Ende des früheren liegt. Bleiben sie
getrennt, liegt dazwischen mindestens ein Schritt. Die Zählung der Zeilen folgt dem Skript, das
jede Zeile beider Folgen genau einmal verbraucht (\in{Abschnitt}[sec:myers]). Der Anfang und die
Länge eines Hunks stimmen daher mit den Positionen in den Folgen überein.

\startinfobox
{\bf Grenzfälle der Hunks.}
\startitemize
\item Ohne Änderung im Skript gibt es keinen Hunk. Zwei gleiche Dateien erscheinen im Befehl gar
nicht. Eine neue oder gelöschte leere Datei erscheint mit ihrem Kopf, aber ohne Hunk.
\item Bei $c = 0$ bilden benachbarte Änderungen ohne unveränderte Zeile dazwischen einen Hunk,
denn ihre Fenster berühren sich. Der Befehl \type{diff} nimmt immer drei Zeilen Kontext.
\item Am Anfang und am Ende des Skripts hat ein Hunk weniger Kontextzeilen.
\item Auf der Seite, die eine neue oder gelöschte Datei nicht hat, ist die Länge des Hunks null.
Die Zeile \type{@@} nennt dann die Zeile davor; die Kommandozeilenreferenz beschreibt das.
\item Der Hunk kennt die Zeilen nicht, nur ihre Positionen. Die Anzeige liest sie aus den Folgen.
\stopitemize
\stopinfobox

\section[sec:diff-command]{Der Diff mit dem Working-Copy-Commit}

Der Befehl \type{diff} vergleicht ohne Argument den Working-Copy-Commit mit seinem Elternteil,
mit einer Kennung diesen Commit mit der Arbeitskopie und mit zwei Kennungen die beiden Commits.
Ist eine Seite der Working-Copy-Commit, hält der Befehl den Stand des Arbeitsverzeichnisses
zuerst in einem Snapshot fest (ADR-010). Erst dann hat die Arbeitskopie eine feste Kennung,
die der Kopf des Diffs nennen kann. Der Befehl liest die Arbeitskopie deshalb nicht nur:
Er kann den Working-Copy-Commit und das Operation Log ändern.

\in{Abbildung}[fig:diff-flow] zeigt den Ablauf. Die Schritte 2 bis 4 sind der Snapshot aus
\in{Kapitel}[chap:phase1], und sie entfallen, wenn der Befehl zwei Commits vergleicht.
Der Vergleich und die Ausgabe lesen nur.

\startplacefigure[title={Ablauf von \type{diff} mit dem Working-Copy-Commit und Absturzpunkten},
    reference=fig:diff-flow,location=here]
\externalfigure[images/diff-flow.pdf][width=0.9\textwidth]
\stopplacefigure

Nach Erfolg steht der Eintrag des Snapshots im Log, wenn sich eine verfolgte Datei geändert hat,
und die Ausgabe zeigt genau den Stand, den der Snapshot festhielt. Ein Fehler vor dem Snapshot,
etwa ein ungültiger Pfad, schreibt nichts. Hat der Snapshot den Zeiger bewegt, die Bewegung aber
nicht sichern können, endet der Befehl mit Exit-Code 3, und die Ausgabe bleibt aus.
Nach einem Absturz im Snapshot gilt, was \in{Abbildung}[fig:oplog-write] zeigt. Der Vergleich
und die Ausgabe ändern nichts, ein Absturz dort lässt also nichts zurück. Der Befehl lässt sich
nach Erfolg, nach einem Fehler und nach einem Absturz wiederholen. Nach einem Erfolg findet ein
zweiter Aufruf den Baum unverändert und schreibt keinen Eintrag.

\in{Abbildung}[fig:diff-timeline] zeigt den Befehl in einer eigenen kurzen Sitzung mit eigenen
Kennungen, nach dem Muster des laufenden Beispiels. Nach dem Commit \quotation{Anfang} ändert
jemand \type{notes.txt} von Hand. Das erste \type{diff} macht einen Snapshot und schreibt den
Eintrag 4. Das zweite, mit \type{--exit-code}, findet den Baum unverändert und schreibt keinen
Eintrag; es endet mit Exit-Code 4, weil sich die Datei unterscheidet. Nach dem Commit der
Änderung vergleicht \type{diff} mit zwei Kennungen die beiden Commits und ändert weder Zeiger
noch Log.

\startplacefigure[title={Änderungen ansehen als Zeitachse},
    reference=fig:diff-timeline,location=here]
\externalfigure[images/diff-timeline.pdf][width=\textwidth]
\stopplacefigure

\startinfobox
{\bf Grenzfälle von \type{diff}.}
\startitemize
\item Ein Pfad, der auf keiner Seite vorkommt, fällt erst nach dem Snapshot auf, weil die Dateien
der Arbeitskopie erst dann feststehen. Der Befehl endet mit Exit-Code 1, und ein Eintrag des
Snapshots bleibt, falls einer entstand.
\item Ein ungültiger Pfad, also ein leerer, ein absoluter oder einer mit \type{..},
schreibt nichts, denn der Befehl prüft ihn vor dem Snapshot.
\item Mit Pfaden beschränkt der Befehl nur die Ausgabe. Der Snapshot umfasst weiter alle
verfolgten Dateien.
\item Ein Pfad im Konflikt (\in{Abschnitt}[sec:conflict-commands]) erscheint wie eine gewöhnliche
Datei, mit dem Text der Markierungen als Inhalt; der Kopf nennt den Pfad mit
\type{(conflict)}.
\stopitemize
\stopinfobox

\section[sec:binary]{Binärdateien}

Ein Zeilenvergleich einer Binärdatei ist sinnlos und kann sehr gross werden. Der Diff erkennt
darum eine Binärdatei und zeigt statt der Zeilen nur, dass sie sich unterscheidet. Die Erkennung
folgt Git: Ein NUL-Byte unter den ersten 8000 Bytes macht die Datei binär. Der Code steht in
\type{diff.rs} als \type{is_binary}. Der Kern wendet sie beim Vergleich der Seiten an
(\type{repository/difference.rs}).

Dieselbe Zusammenfassung gilt für eine Textdatei, die zu gross für den Zeilenvergleich ist.
Die Grenze liegt bei 10 MiB je Seite (\type{DIFF_SIZE_LIMIT}, Prüfung \type{is_too_large}).
Die Begründung steht in den Entscheidungen oben: Zeit und Ausgabe wachsen mit der Datei,
die Grenze hält beide in Schranken.

\startitemize[n]
\item Für jede Datei, die sich unterscheidet, liest der Kern die Inhalte beider Seiten. Eine
Seite, auf der die Datei fehlt, gilt als leer.
\item Ist eine der beiden Fassungen binär, ist die ganze Datei binär. Der Kern vergleicht dann
keine Zeilen und bildet keine Hunks.
\item Ist keine Fassung binär, aber eine grösser als 10 MiB, ist die Datei zu gross. Der Kern
vergleicht auch dann keine Zeilen und bildet keine Hunks.
\item Sonst vergleicht der Kern die Zeilen und bildet die Hunks.
\item Die Anzeige schreibt für eine binäre Datei die Zeile \type{Binary files ... and ... differ}
mit den beiden Seiten wie im Kopf. Die Oberfläche zeigt \type{Binary file differs}.
\item Für eine zu grosse Datei schreibt die Anzeige die Zeile
\type{Files ... and ... differ and are too large to compare}, ebenfalls mit den beiden Seiten.
Die Oberfläche zeigt \type{File too large to compare}.
\stopitemize

\subsection{Definition und Komplexität}

Sei $c = (c_0, \dots, c_{L-1})$ der Inhalt einer Datei der Länge $L$. Die Datei ist binär, genau
dann wenn es ein $i < \min(L, 8000)$ mit $c_i = 0$ gibt. Die Prüfung liest höchstens 8000 Bytes
und braucht daher konstante Zeit und konstanten Speicher. Das Lesen des Inhalts, das ihr
vorausgeht, kostet $O(L)$.

\subsection{Garantien}

Die Prüfung ist eine reine Funktion des Inhalts. Der Kern entscheidet einmal, und die Textausgabe
und die Oberfläche übernehmen das Ergebnis. Dieselbe Datei erscheint also in beiden gleich. Die
Prüfung ändert nichts im Repository.
Die Prüfung der Grösse vergleicht nur die Länge und kostet nach dem Lesen $O(1)$.

\startinfobox
{\bf Grenzfälle der Erkennung.}
\startitemize
\item Die Erkennung ist eine Heuristik wie in Git. Eine Textdatei in UTF-16 enthält NUL-Bytes und
gilt als binär. Eine Binärdatei ohne NUL-Byte in den ersten 8000 Bytes gilt als Text.
\item Steht das erste NUL-Byte erst ab dem 8001. Byte, ist die Datei Text.
\item Eine leere Datei ist keine Binärdatei.
\item Wird aus einer Textdatei eine Binärdatei oder umgekehrt, gilt die ganze Datei als binär.
\item Eine neue oder gelöschte Binärdatei erscheint mit \type{/dev/null} auf der Seite,
auf der sie fehlt.
\item Eine Datei, die binär ist, bleibt binär, auch wenn sie über der Grössengrenze liegt.
Die Grenze gilt nur für Text.
\item Der Kern liest eine grosse Datei ganz, ehe er ihre Grösse prüft. Die Grenze schützt also
vor der Zeit und der Ausgabe des Vergleichs, nicht vor dem Lesen und nicht vor dem Speicher,
den beide Seiten dabei belegen.
\stopitemize
\stopinfobox

\section[sec:pager]{Farbe und Pager}

Der Text des Befehls \type{diff} ist die Referenz für Skripte und Tests. Auf einem Terminal soll er
auch bequem zu lesen sein. Darum färbt der Befehl die Zeilen wie Git und gibt den Text an einen
Pager weiter, wenn die Standardausgabe ein Terminal ist (\in{Abbildung}[fig:diff-output]). Der
Code steht in \type{cli.rs} als \type{diff_text} und \type{write_paged}.

\startplacefigure[title={Wege der Ausgabe von \type{diff}},
    reference=fig:diff-output,location=here]
\externalfigure[images/diff-output.pdf][width=0.9\textwidth]
\stopplacefigure

\startitemize[n]
\item Der Kern berechnet den Diff. Ist eine Seite der Working-Copy-Commit, speichert er vorher
die verfolgten Dateien als Snapshot (\in{Abschnitt}[sec:diff-command]). Danach ändert er nichts
mehr.
\item Der Befehl baut den Text. Mit Farbe sind die Zeilen \type{diff}, \type{---} und
\type{+++} fett, die Zeilen \type{@@} cyan, entfernte Zeilen rot und eingefügte grün. Entfernte
Zeilen, die an anderer Stelle wieder auftauchen, und die dort eingefügten Zeilen haben eine
eigene Farbe (\in{Abschnitt}[sec:moved]).
Die Farbe endet vor dem Zeilenende, damit sie nicht in die nächste Zeile läuft.
\item Auf einem Terminal unter Unix startet er eine Shell, die den Befehl aus \type{PAGER}
ausführt, ersatzweise \type{less}. Ist \type{LESS} nicht gesetzt, bekommt der Pager
\type{LESS=FRX}. Der Befehl schreibt den Text in den Eingang des Pagers und wartet, bis der Pager
endet.
\item Sonst schreibt er den Text ohne Farbe auf die Standardausgabe.
\stopitemize

\subsection{Garantien}

Der Text ohne Farbe ist der Text mit Farbe ohne die Escape-Folgen. Ein Skript, das eine Pipe liest,
bekommt also immer reinen Text, ohne Farbe und ohne Pager. Der Diff geht nicht verloren: Lässt
sich der Pager nicht starten, oder meldet die Shell den Exit-Code 126 oder 127 (ein Befehl, der
nicht ausführbar ist oder fehlt), schreibt der Befehl den Text ungefiltert, also mit Farbe, auf die
Standardausgabe und einen Hinweis auf stderr. Ein Pager, der früher endet, oder ein anderer
Exit-Code des Pagers ist kein Fehler, und eine geschlossene Pipe beendet die Ausgabe leise.
Der Exit-Code des Befehls ist das Ergebnis des Diffs, also 0 oder mit \type{--exit-code} bei
Unterschieden 4, nie der des Pagers. Ausnahmen sind ein Fehler beim Schreiben (1) und ein Signal,
das den Befehl beendet (130, 143 oder 129).

Ein Absturz in diesem Ablauf hinterlässt nichts, das zu reparieren wäre. Der Snapshot ist
gespeichert, ehe der Pager startet, und der Rest schreibt nur auf das Terminal. Der Befehl lässt
sich danach gefahrlos wiederholen.

Das Verhalten bei Signalen folgt aus einer Eigenschaft des Terminals. Ctrl-C geht an alle Prozesse
im Vordergrund, also an den Pager und an den Befehl. Endet der Befehl zuerst, übernimmt die Shell
das Terminal, während der Pager es noch liest, und beide streiten um die Eingabe. Der Befehl
wartet darum, bis der Pager endet, und endet erst dann mit seinem eigenen Exit-Code. Ein Signal
zum Beenden, das nur den Befehl trifft, etwa ein \type{kill} auf seine Prozessnummer, sendet er
an die Shell des Pagers, und er endet mit 143, bei \type{SIGHUP} mit 129.

\startinfobox
{\bf Grenzfälle von Farbe und Pager.}
\startitemize
\item Der Pager ist der Befehl in \type{PAGER}. Eine leere oder fehlende Variable gilt als nicht
gesetzt. \type{GIT_PAGER} und die Einstellung \type{core.pager} von Git beachtet der Befehl nicht.
\item Der Buchstabe \type{R} in \type{LESS=FRX} reicht die Farben durch, \type{F} beendet
\type{less} bei einer einzigen Seite, und \type{X} lässt den Bildschirm des Terminals stehen.
\item Ist \type{PAGER} eine längere Befehlszeile, bleibt die Shell bestehen. Trifft ein
\type{SIGTERM} nur den Befehl, endet die Shell, aber ihre Befehle laufen weiter,
bis das Terminal schliesst.
\item Unter Windows gibt es in Phase 2 immer reinen Text, denn die Konsole braucht eine
Verarbeitung der Escape-Folgen und eine Shell für den Pager.
\item Eine Option oder Variable, die die Farbe erzwingt oder abschaltet, gibt es in Phase 2
nicht.
\stopitemize
\stopinfobox

\section[sec:moved]{Verschobene Zeilen}

Ein Block, der von einer Stelle an eine andere wandert, erscheint im Diff als Löschung und als
Einfügung (die Entscheidung steht in \type{docs/research/moved-lines-heuristics.md}).
Das bleibt so, denn die Hunks sind die Grundlage für
Skripte, Tests und den Merge-Pfad. Auf dem Terminal markiert die Anzeige jedoch, welche Zeilen
verschoben wurden, wie \type{git diff --color-moved}. Die Markierung ist reine Anzeige: Sie
ändert weder die Hunks noch den Text ohne Farbe noch den Exit-Code.

\subsection{Erkennung}

Die Funktion \type{moved_lines} im Modul \type{moved} des Kerns liest einen fertigen Diff und
ändert ihn nicht. Sie arbeitet auf allen Textdateien des Diffs zugleich, ein Block kann also
in eine andere Datei wandern.

\startitemize[n]
\item Sie sammelt die gelöschten und die eingefügten Zeilen aller Hunks in der Reihenfolge des
Diffs. Ein Block ist eine Folge von Zeilen, die auf ihrer Seite unmittelbar aufeinander folgen.
\item Zu jeder gelöschten Zeile sucht sie den längsten Block eingefügter Zeilen, der mit ihr
beginnt und Zeile für Zeile gleich ist, Byte für Byte und samt Zeilenende.
Bei gleicher Länge gilt der erste.
\item Der Block zählt nur, wenn seine Zeilen zusammen mindestens 20 alphanumerische Zeichen
haben, dieselbe Zahl wie in Git, gezählt über die Unicode-Zeichen.
Schliessende Klammern und Leerzeilen allein gelten deshalb nicht als Verschiebung.
\item Eine Zeile gehört zu höchstens einer Verschiebung. Pro gelöschter Zeile probiert die
Funktion höchstens 256 Stellen auf der neuen Seite, und insgesamt hat sie ein Budget von
20 Millionen Zeilenvergleichen pro Diff (der letzte Block kann es knapp überschreiten).
Eine oft wiederholte Zeile dehnt die Laufzeit so nicht;
ist das Budget aufgebraucht, behält die Funktion die bis dahin gefundenen Verschiebungen.
\stopitemize

Der Ablauf ist deterministisch: Derselbe Diff liefert dieselben Markierungen.
Eine Zeile, die beim Verschieben geändert wurde, unterbricht den Block, und die geänderte
Zeile bleibt eine gewöhnliche Löschung und Einfügung.
Die Teile davor und danach zählen einzeln gegen die Schwelle.
Binärdateien und Dateien über der Grössengrenze nehmen nicht teil.

\subsection{Farben}

Die Farben sind die von Git. Eine entfernte Zeile, die verschoben wurde, ist fett magenta,
und eine eingefügte Zeile, die von anderswo kommt, ist fett cyan. Das Zeichen am Zeilenanfang
bleibt \type{-} oder \type{+}. Der Text ohne Farbe zeigt deshalb keine Verschiebung, und eine
Pipe oder eine Datei bekommt weiter reinen Text. Die Oberfläche von \type{diff --tui} verwendet
dieselben Farben für ihre Zeilen (\in{Abschnitt}[sec:tui]).

\section[sec:tui]{Die Oberfläche im Terminal}

\type{diff --tui} zeigt denselben Diff in einer schreibgeschützten Oberfläche: links die Dateien,
rechts die Zeilen aller Dateien in einem durchgehenden Strom, unten eine Zeile mit den Tasten.
Jede Datei beginnt im Strom mit einer Kopfzeile, die den Pfad und die Zahl der eingefügten und
entfernten Zeilen nennt (\type{+n -m}). Die Liste links markiert die Datei, die oben im Bild steht.
Die Bibliothek ist ratatui mit dem Backend crossterm (ADR-016); die Syntaxhervorhebung
kommt von tree-sitter (ADR-021). Die Crate \type{stemma-tui}
bekommt den fertigen Diff von der Kommandozeile und berechnet ihn nicht neu; sie braucht seine
Typen und \type{moved_lines}, um verschobene Zeilen zu färben (\in{Abschnitt}[sec:moved]).
Der Kern hängt nicht von ihr ab, und der Text von \type{diff} bleibt die Bezugsgrösse für Skripte
und Tests.
Der Zustand ist ein einfacher Wert, der sich durch Eingaben ändert, und die
Darstellung zeichnet ihn in einen beliebigen Bildschirm. Tests brauchen daher kein Terminal. Nur
die Schleife der Oberfläche (\type{run}) berührt das echte Terminal.

\subsection{Der Strom}

Der Strom ist die Liste der Zeilen, die das rechte Feld zeigt. Er enthält nacheinander für jede
Datei eine Kopfzeile, danach bei einer Textdatei für jeden Hunk die Zeile \type{@@} und seine
Zeilen, bei einer Binärdatei oder einer zu grossen Datei eine Zeile mit der Meldung. Zwischen
zwei Dateien steht eine Leerzeile; sie gehört zur Datei darüber. Eine Zeile des Stroms heisst
im Folgenden kurz Zeile, eine Zeile einer Datei Dateizeile.

Der Strom hängt von der Breite des Feldes, vom Layout und vom Umbruch ab. Die Oberfläche baut
ihn einmal für diese drei Werte und nicht für jedes Bild neu; sie baut ihn erst wieder, wenn sich
einer von ihnen ändert. Dabei zählt die Breite in Spalten des Terminals, nicht in Bytes: Ein
breites Zeichen wie ein chinesisches nimmt zwei Spalten ein. Ein Zeichen, das über den Rand
reichen würde, kommt in die nächste Zeile; ohne Umbruch wird die Zeile abgeschnitten, und die
letzte Spalte zeigt \type{…}, damit der Schnitt zu sehen ist. Ein Zeichen, das breiter ist als
die ganze Spalte, erscheint als \type{…}.

\startitemize[n]
\item {\bf Unified.} Eine Dateizeile ist eine Zeile des Stroms. Vor der Marke (Leerzeichen,
\type{-} oder \type{+}) stehen zwei Spalten mit Zeilennummern, links die der alten und
daneben die der neuen Seite. Eine entfernte Zeile hat nur die alte, eine eingefügte nur die
neue Nummer; eine Kontextzeile hat beide. Die Spalten sind so breit wie die Nummer der
längsten Datei.
\item {\bf Split.} Links steht die alte, rechts die neue Seite, getrennt durch einen senkrechten
Strich. Jede Seite hat ihre eigene Nummer und Marke. Eine Kontextzeile steht auf beiden Seiten.
Eine zusammenhängende Gruppe aus Löschungen und Einfügungen wird Zeile für Zeile gepaart: die
erste Löschung mit der ersten Einfügung, die zweite mit der zweiten, und die Seite mit weniger
Zeilen bleibt darunter leer. Ein Paar belegt eine Zeile des Stroms; ohne Umbruch hat der Strom
daher im Layout Split höchstens so viele Zeilen wie im Layout Unified.
\item {\bf Auto.} Das Layout ist Split, wenn das Feld mindestens 120 Spalten innerhalb seines
Rahmens hat, und sonst Unified. Jede Hälfte hat dann 59 oder 60 Spalten, von denen die Nummer
und die Marke einige nehmen. Die Liste der Dateien nimmt ein Viertel des Terminals, aber nicht
mehr als 40 Spalten, und der Rahmen des Feldes zwei; ein Terminal von 162 Spalten und mehr
teilt also, eines von 161 Spalten nicht.
Die Tasten \type{1}, \type{2} und \type{0} wählen Unified, Split und Auto.
\item {\bf Umbruch.} Ohne Umbruch (der Anfangszustand) schneidet die Oberfläche eine Zeile am
Rand der Spalte ab und zeigt \type{…} in der letzten Spalte. Mit Umbruch (Taste \type{w}) geht
sie in der nächsten Zeile des Stroms weiter, deren Nummernspalten und Marke leer sind. Die Zeile
\type{\ No newline at end of file} wird in beiden Layouts ebenso umgebrochen oder abgeschnitten.
\stopitemize

Entfernte Zeilen haben einen roten, eingefügte einen grünen Hintergrund, der bis an den Rand der
Spalte reicht. Den Text der Zeilen färbt die Syntaxhervorhebung
(\in{Abschnitt}[sec:tui-highlight]).
Verschobene Zeilen (\in{Abschnitt}[sec:moved]) behalten die Vordergrundfarben
magenta und cyan ohne Hintergrund und bleiben so von gewöhnlichen Änderungen zu unterscheiden.
Ein Tabulator erscheint als vier Leerzeichen. Ein anderes Steuerzeichen erscheint in
Caret-Schreibweise, etwa \type{^[} für Escape, \type{^M} für ein CR mitten in der Zeile und
\type{^?} für DEL, ein Zeichen des Bereichs C1 als \type{^@}; so erreicht keine Folge aus der
Datei das Terminal. Das CR eines Zeilenendes \type{\r\n} gehört zum Zeilenende und entfällt.
Ein Pfad, der länger ist als das Feld, wird in der Kopfzeile abgeschnitten und endet auf
\type{…}, damit die Zahlen der Kopfzeile sichtbar bleiben.

\subsection[sec:tui-highlight]{Syntaxhervorhebung und Themen}

Die Oberfläche färbt den Text der Zeilen nach der Syntax der Datei (ADR-021). Sie wählt die
Grammatik nach der Endung des Pfads: \type{.rs} (Rust), \type{.toml}, \type{.sh} und \type{.bash},
\type{.json} sowie \type{.md} (nur die Blockstruktur von Markdown, also Überschriften und
Codeblöcke). Eine Datei mit einer anderen Endung, und ein Text, den die Grammatik nicht lesen
kann, behalten die Farben ohne Hervorhebung; ein Fehler entsteht dabei nie.

Die Hervorhebung liest den ganzen Text einer Seite auf einmal, die alte und die neue Seite je
für sich, und nicht Zeile für Zeile. Ein Block-Kommentar oder ein mehrzeiliger String ist so in
jeder seiner Zeilen gefärbt, auch wenn der Diff nur eine davon zeigt. Der Parser (tree-sitter)
ist fehlertolerant: Ein halber Satz oder eine Datei mit einem Syntaxfehler liefert trotzdem
Farben für die Teile, die er erkennt. Die Bereiche werden beim ersten Bau des Stroms bestimmt und
danach behalten, auch wenn das Thema wechselt.

Die Hervorhebung ändert nur die Vordergrundfarbe der Zeichen. Der Hintergrund einer entfernten
oder eingefügten Zeile, der Text, die Breite und der Umbruch bleiben, wie sie ohne sie wären.
Verschobene Zeilen (\in{Abschnitt}[sec:moved]) werden nicht hervorgehoben, denn ihre Farbe
sagt, dass sie verschoben sind.

Ein Thema legt die Farben fest. {\bf Dunkel} (der Anfangszustand) ist das Aussehen von früher für
ein Terminal mit dunklem Hintergrund, mit Hervorhebung. {\bf Hell} hat dunklen Text auf hellen
Hintergründen für ein helles Terminal. {\bf Schlicht} behält die Farben von Dunkel und schaltet
die Hervorhebung aus. Die Taste \type{t} wechselt zum nächsten Thema in dieser Reihenfolge, und
die Titelzeile des Feldes nennt ein Thema, das nicht Dunkel ist. Ein Thema hat keine
Einstellung ausserhalb der Oberfläche; es gilt, bis man sie beendet.

\subsection{Definition}

Der Zustand ist $S = (f, h, s, m, w, c, t)$ (im Code \type{State}). Dabei ist $f$ die Nummer der
Datei im Bild, nämlich der Datei, zu der die oberste Zeile des Feldes gehört. Die Nummer $h$ ist
die des gewählten Hunks in dieser Datei, und $s$ ist die Zahl der Zeilen des Stroms, die oben aus
dem Bild gerollt sind. Das Layout $m$ ist eines von Auto, Unified und Split, der Umbruch
$w$ ist an oder aus, $c$ ist die Breite des Terminals in Spalten, die der Zustand durch die
Eingabe \type{Resize} erfährt, und $t$ ist das Thema der Farben (Dunkel, Hell oder Schlicht,
\in{Abschnitt}[sec:tui-highlight]). Der Diff $D$ ist fest und wird bei jedem Aufruf mitgegeben.
Er hat $F$ Dateien. Die Datei $f$ hat $H(f)$ Hunks. $H(f) = 0$ gilt für eine binäre Datei und
für eine neue oder gelöschte leere Datei. Der Strom hat $R$ Zeilen; er hängt von $c$, $m$
und $w$ ab. Mit $\mathrm{kopf}(f)$ sei der Index der Kopfzeile der Datei $f$ im Strom bezeichnet,
mit $\mathrm{start}(f, h)$ der Index der Zeile \type{@@} des Hunks $h$, und
$\mathrm{start}(f, 0) = \mathrm{kopf}(f)$ für eine Datei ohne Hunk. Die Funktion
$\mathrm{datei}(s)$ liefert die Datei, zu der die Zeile $s$ gehört. Der Anker der Auswahl ist
$\mathrm{anker}(f, h) = \mathrm{kopf}(f)$ für $h = 0$ und $\mathrm{start}(f, h)$ sonst: die Zeile,
die oben stehen soll, wenn die Auswahl $(f, h)$ ist.

Eine Eingabe (im Code \type{Input}) ist eine der Tasten oder eine der Mausaktionen aus
\in{Tabelle}[tab:tui-transitions], oder die Meldung \type{Resize}, die die Schleife schickt,
wenn sich die Breite des Terminals ändert.
Die Maus ist wie die Tasten ein einfacher Wert: Ein Klick trägt die Spalte und die Zeile der
Zelle. Die Oberfläche merkt sich, wo das letzte Bild die Namen der Dateien gezeichnet hat
($\mathrm{liste}(x, y)$ liefert die Datei, deren Name in der Zelle $(x, y)$ steht), und bildet
den Klick damit auf eine Datei ab; so brauchen auch die Tests der Maus kein Terminal.
Die Tastatur bleibt vollständig, die Maus ist nur eine Abkürzung.
Die Funktion $\mathrm{apply}$ bildet einen Zustand und eine Eingabe auf den nächsten
Zustand ab.
Sie liefert zugleich, ob die Oberfläche weiterläuft: Nur das Beenden liefert \type{false}.
Die Tabelle nennt nur die Teile des Zustands, die sich ändern; die übrigen bleiben.

\startplacetable[title={Übergänge des Zustands $(f, h, s, m, w, c, t)$},
  reference=tab:tui-transitions,location=here]
\startxtable[frame=off,option=stretch]
\startxrow[topframe=on,bottomframe=on]
\startxcell {\bf Eingabe} \stopxcell
\startxcell {\bf Tasten} \stopxcell
\startxcell {\bf Bedingung} \stopxcell
\startxcell {\bf Neuer Zustand} \stopxcell
\stopxrow
\startxrow
\startxcell Beenden \stopxcell
\startxcell \type{q}, Esc, Ctrl-C \stopxcell
\startxcell immer \stopxcell
\startxcell unverändert, die Oberfläche endet \stopxcell
\stopxrow
\startxrow
\startxcell Nächste Datei \stopxcell
\startxcell \type{J}, Tab \stopxcell
\startxcell $f + 1 < F$ \stopxcell
\startxcell $f := f+1$, $h := 0$, $s := \mathrm{kopf}(f+1)$ \stopxcell
\stopxrow
\startxrow
\startxcell Vorige Datei \stopxcell
\startxcell \type{K}, Shift-Tab \stopxcell
\startxcell $f > 0$ \stopxcell
\startxcell $f := f-1$, $h := 0$, $s := \mathrm{kopf}(f-1)$ \stopxcell
\stopxrow
\startxrow
\startxcell Nächster Hunk \stopxcell
\startxcell \type{n} \stopxcell
\startxcell $h + 1 < H(f)$ \stopxcell
\startxcell $h := h+1$, $s := \mathrm{start}(f, h+1)$ \stopxcell
\stopxrow
\startxrow
\startxcell \stopxcell
\startxcell \stopxcell
\startxcell $h + 1 \ge H(f)$ und $f + 1 < F$ \stopxcell
\startxcell $f := f+1$, $h := 0$, $s := \mathrm{kopf}(f+1)$ \stopxcell
\stopxrow
\startxrow
\startxcell Voriger Hunk \stopxcell
\startxcell \type{p} \stopxcell
\startxcell $h > 0$ \stopxcell
\startxcell $h := h-1$, $s := \mathrm{start}(f, h-1)$ \stopxcell
\stopxrow
\startxrow
\startxcell \stopxcell
\startxcell \stopxcell
\startxcell $h = 0$ und $f > 0$ \stopxcell
\startxcell $f := f-1$, $h := h'$, $s := \mathrm{start}(f-1, h')$
mit $h' = \max(H(f-1) - 1, 0)$ \stopxcell
\stopxrow
\startxrow
\startxcell Abwärts blättern um $k$ \stopxcell
\startxcell \type{j}, Pfeil ab ($k = 1$); Leertaste, Bild ab ($k = 20$) \stopxcell
\startxcell immer \stopxcell
\startxcell $s := \min(s + k,\ \max(R - 1, 0))$; wechselt dadurch die Datei, ist
$f := \mathrm{datei}(s)$ und $h := 0$ \stopxcell
\stopxrow
\startxrow
\startxcell Aufwärts blättern um $k$ \stopxcell
\startxcell \type{k}, Pfeil auf ($k = 1$); \type{b}, Bild auf ($k = 20$) \stopxcell
\startxcell immer \stopxcell
\startxcell $s := \max(s - k,\ 0)$; wechselt dadurch die Datei, ist
$f := \mathrm{datei}(s)$ und $h := 0$ \stopxcell
\stopxrow
\startxrow
\startxcell Mausrad $k = 3$ \stopxcell
\startxcell Rad ab, Rad auf \stopxcell
\startxcell immer \stopxcell
\startxcell wie Abwärts und Aufwärts blättern um $k$ \stopxcell
\stopxrow
\startxrow
\startxcell Datei anklicken \stopxcell
\startxcell linke Maustaste auf den Namen einer Datei in der Liste \stopxcell
\startxcell $\mathrm{liste}(x, y) = g$ \stopxcell
\startxcell $f := g$, $h := 0$, $s := \mathrm{kopf}(g)$, wie bei den Tasten für Dateien \stopxcell
\stopxrow
\startxrow
\startxcell \stopxcell
\startxcell \stopxcell
\startxcell die Zelle liegt ausserhalb der Namen, oder es gab noch kein Bild \stopxcell
\startxcell unverändert \stopxcell
\stopxrow
\startxrow
\startxcell Layout wählen \stopxcell
\startxcell \type{1} Unified, \type{2} Split, \type{0} Auto \stopxcell
\startxcell immer \stopxcell
\startxcell $m := $ gewähltes Layout; der Strom wird neu gebaut und
$s := \min(\mathrm{anker}(f, h),\ \max(R' - 1, 0))$ \stopxcell
\stopxrow
\startxrow
\startxcell Umbruch umschalten \stopxcell
\startxcell \type{w} \stopxcell
\startxcell immer \stopxcell
\startxcell $w := \neg w$; der Strom wird neu gebaut und $s$ wie beim Layout \stopxcell
\stopxrow
\startxrow
\startxcell Thema wechseln \stopxcell
\startxcell \type{t} \stopxcell
\startxcell immer \stopxcell
\startxcell $t :=$ das nächste Thema (Dunkel, Hell, Schlicht, Dunkel); die Zeilen behalten ihre
Zahl, darum bleibt $s$ \stopxcell
\stopxrow
\startxrow[bottomframe=on]
\startxcell Breite melden \stopxcell
\startxcell keine Taste, \type{Resize} \stopxcell
\startxcell immer \stopxcell
\startxcell $c := $ neue Breite; der Strom wird neu gebaut und $s$ wie beim Layout \stopxcell
\stopxrow
\stopxtable
\stopplacetable

Dabei ist $R'$ die Zahl der Zeilen des neuen Stroms. Trifft keine Bedingung einer Zeile zu, etwa
\type{J} auf der letzten Datei, bleibt der Zustand unverändert. Der Ausgangszustand ist
$(0, 0, 0, \mathrm{Auto}, \mathrm{aus}, 80, \mathrm{Dunkel})$, bis die Schleife die wahre
Breite meldet.
Der Ablauf ist in \in{Abbildung}[fig:tui-states] gezeichnet.

\startplacefigure[title={Zustände der Oberfläche von \type{diff --tui}},
    reference=fig:tui-states,location=here]
\externalfigure[images/tui-states.pdf][width=\textwidth]
\stopplacefigure

\subsection{Invariante}

Für jeden Zustand, den die Übergänge aus dem Ausgangszustand erreichen, gilt: $f < F$
(oder $f = 0$ bei einem leeren Diff), $h < \max(H(f), 1)$ und $s \le \max(R - 1, 0)$.
Ausserdem gehört die oberste Zeile zur Datei im Bild: $f = \mathrm{datei}(s)$, wenn der Diff
Dateien hat. Die Datei wechselt nur mit den Tasten für Dateien und Hunks und dadurch, dass das
Blättern eine andere Datei an den oberen Rand bringt; jeder Wechsel setzt $h$ neu. Blättern
innerhalb einer Datei ändert $f$ und $h$ nicht. Darum gehen \type{n} und \type{p} vom zuletzt
gewählten Hunk aus und nicht von dem, den das Blättern ins Bild gebracht hat.

\subsection{Komplexität}

Sei $T$ die Zahl der Zeichen in den Hunks aller Dateien. Den Strom zu bauen kostet
$O(T + R \cdot c)$ für jede Änderung von $c$, $m$ oder $w$, und er braucht $O(R \cdot c)$
Speicher, weil jede Zeile bis an den Rand aufgefüllt ist. Die Zeilen hängen damit nicht mehr vom
Zeichnen ab: Die Wahl einer Datei oder eines Hunks, das Blättern und das Zeichnen lesen den
fertigen Strom. Eine Datei oder einen Hunk zu wählen kostet $O(1)$, und
$\mathrm{datei}(s)$ kostet $O(\log F)$ durch binäre Suche. Das Zeichnen kostet $O(h_t + F)$ mit
der Höhe $h_t$ des Terminals: Es kopiert nur die Zeilen im Bild und listet alle $F$ Dateien.
Ausser dem Zustand hält die Oberfläche den Strom, aber keine Kopie des Diffs selbst.

\subsection{Korrektheit}

Die Invarianten folgen aus der Tabelle. Jede Zeile hält die Schranken ein: Der Wechsel zur
nächsten Datei prüft $f + 1 < F$, der zum nächsten Hunk $h + 1 < H(f)$,
und $\mathrm{kopf}$ und $\mathrm{start}$ liefern den Index einer Zeile des Stroms.
Blättern nach unten begrenzt $s$ durch die letzte Zeile, nach oben durch null, und $f$ wird
danach aus $s$ bestimmt. Beim Wechsel des Layouts, des Umbruchs oder der Breite fällt die Zahl der
Zeilen anders aus; $s$ wird deshalb auf die Zeile des Ankers gesetzt und durch die neue letzte
Zeile begrenzt. Die Auswahl bleibt dabei erhalten, nur ihre Zeile im Strom ändert sich. Die
Auswahl stand an einer Zeile, die es auch im neuen Strom gibt, denn jede Datei und jeder Hunk
hat in jedem Layout genau eine Kopfzeile und eine Zeile \type{@@}. Die Funktion $\mathrm{apply}$
ist total: Jede Eingabe hat ein Ergebnis, und kein Zustand lässt sie scheitern. Ein Zustand, der
zu einem anderen Diff gehört, zeigt ein leeres Feld, löst aber keinen Fehler aus. Weil
$\mathrm{apply}$ nur den Zustand, den Diff und den Strom liest, prüfen die Tests sie ohne Terminal.
Dasselbe gilt für ein sehr kleines Terminal: Eine Spalte kann die Breite null haben, und sie
zeigt dann nichts, auch kein \type{…}. Was nicht hineinpasst, schneidet die Oberfläche oder
ratatui ab, ohne zu scheitern.

\subsection{Garantien}

Die Oberfläche liest nur. Der Befehl hat den Snapshot vor dem Start gemacht, wie bei der
Textausgabe, und danach schreibt er nichts mehr ins Repository. Die Oberfläche gibt das Terminal
bei jedem Ende zurück: beim Beenden, bei einem Fehler und bei einer Panik, für die ratatui einen
Haken setzt. Die Maus schaltet die Oberfläche mit sich ein und bei jedem dieser Enden wieder aus;
den Haken für die Panik erweitert sie darum, denn der von ratatui kennt die Maus nicht.
Text lässt sich im Terminal weiter markieren, solange man die Taste hält, die das Terminal dafür
vorsieht (meist Umschalt, in iTerm2 die Wahltaste); die Hilfe von \type{diff --tui} sagt es.
Gelingt die Rückgabe nach einem Fehler nicht, nennt die Meldung beide Fehler.
Ein Handler für Signale läuft auf einem anderen Thread und fasst das Terminal nicht an. Er setzt
eine Marke, die die Oberfläche mindestens zehnmal in der Sekunde und nach jedem Ereignis liest,
auch nach einem der Maus. Die
Oberfläche beendet sich dann und gibt das Terminal zurück, und erst danach endet der Befehl mit
dem Exit-Code des Signals. Der Handler wartet darauf höchstens zwei Sekunden.

\startinfobox
{\bf Grenzfälle der Oberfläche.}
\startitemize
\item Ohne Unterschiede startet die Oberfläche nicht, und der Befehl endet ohne Ausgabe.
\item Ist die Standardausgabe kein Terminal oder läuft der Befehl unter Windows, endet er mit
Exit-Code 1, ehe er etwas liest oder speichert.
\item Ein Tabulator erscheint als vier Leerzeichen, nicht ausgerichtet auf Tabulatorstopps.
Ein anderes Steuerzeichen erscheint in Caret-Schreibweise (\type{^[}, \type{^M}, \type{^?}),
das CR vor dem Zeilenende erscheint nicht.
\item Eine Zeile ohne Zeilenende erhält die Zeile \type{\ No newline at end of file}.
\item Eine binäre Datei zeigt die Meldung \type{Binary file differs}, eine zu grosse Datei die
Meldung \type{File too large to compare}; beide haben keinen Hunk.
Eine leere neue oder gelöschte Datei zeigt ein leeres Feld.
\item Ein \type{SIGKILL} lässt sich nicht abfangen, und bei einem Signal ohne Handler bleibt das
Terminal im Raw-Modus. \type{reset} stellt es wieder her.
\item Ctrl-C kommt im Raw-Modus als Taste an und wirkt wie \type{q}: Der Befehl endet mit 0, mit
\type{--exit-code} und Unterschieden mit 4, nicht mit 130.
\stopitemize
\stopinfobox

\section[sec:diff-workflows]{Änderungen ansehen: die Arbeitsabläufe}

Der Diff ist ein lesender Befehl mit einer Ausnahme: Ist eine Seite der Working-Copy-Commit,
macht er zuerst einen Snapshot der verfolgten Dateien (ADR-010), wie \type{status}.
Der Snapshot schreibt einen Eintrag nur, wenn sich eine verfolgte Datei geändert hat.
Vergleicht der Aufruf zwei Commits, liest er nur Objekte und schreibt gar nichts.
Einen Akteur verlangt er in beiden Fällen, auch wenn nichts zu schreiben ist.
Die folgenden Abläufe greifen Stellen der Sitzung aus \in{Kapitel}[chap:phase1] auf.
Die Ausgaben stammen aus einem echten Lauf. Die Zeit gehört zum Commit, ein neuer Lauf
erzeugt also andere Hashes; die Beispiele zeigen deshalb die gekürzten IDs der Sitzung aus
\in{Kapitel}[chap:phase1]. Ein Aufruf braucht die volle Kennung.
Die Synopsis, die Exit-Codes und die Fehler jeder Option stehen in der Befehlsreferenz.

\subsection{Änderungen vor dem Commit ansehen}

Wer \type{README.md} geändert hat, will vor dem Commit lesen, was er festhält.
Der Ablauf ist der des zweiten Commits in \in{Kapitel}[chap:phase1]:

\starttyping
$ echo "Hallo Welt" > README.md
$ stemma diff
diff README.md
--- b3:2037…df79 README.md
+++ working-copy b3:7694…d054 README.md
@@ -1 +1 @@
-Hallo
+Hallo Welt
\stoptyping

Die linke Seite ist der Elternteil des Working-Copy-Commits, der Commit \quotation{Erste Fassung};
die rechte Seite ist der Working-Copy-Commit selbst, genannt mit der ID seines Snapshots.
Das Arbeitsverzeichnis bleibt unberührt, der Befehl schreibt keine Datei hinein.
Der Snapshot speichert die geänderte \type{README.md} und bewegt den Zeiger des
Working-Copy-Commits von \type{b3:e6a3…} auf \type{b3:7694…}; \type{op log} zeigt dazu den
Eintrag 4, \type{snapshot}. Der Branch \type{main} bleibt stehen.
Ein zweiter Aufruf liest denselben Stand, findet nichts zu speichern und schreibt keinen Eintrag.
Deshalb schreiben auch das folgende \type{status} und der Snapshot im \type{commit}
nichts mehr: Die Änderung ist schon im Working-Copy-Commit.

Eine neue Datei erscheint erst, wenn \type{track} sie aufgenommen hat, denn der Diff
vergleicht nur verfolgte Dateien. Nennt man ihren Pfad vorher, endet der Befehl mit einem Fehler
und nennt \type{track} als Abhilfe. Der Fehler fällt erst nach dem Snapshot auf, der Eintrag
steht dann schon im Log, und die Meldung nennt ihn.

\subsection{Zwei Commits vergleichen}

Nach dem zweiten Commit ist der Working-Copy-Commit leer, ein Diff ohne Argumente also auch.
Was die beiden Commits der Sitzung unterscheidet, zeigt der Aufruf mit beiden Kennungen:

\starttyping
$ stemma diff b3:2037…df79 b3:9d8a…4c77
diff README.md
--- b3:2037…df79 README.md
+++ b3:9d8a…4c77 README.md
@@ -1 +1 @@
-Hallo
+Hallo Welt
\stoptyping

Beide Seiten sind Commits, deshalb gibt es keinen Snapshot.
Das Arbeitsverzeichnis, der Working-Copy-Commit, die Zeiger und das Log bleiben, wie sie sind;
der Befehl kann jederzeit wiederholt werden.
Mit nur einer Kennung vergleicht der Befehl diesen Commit mit dem Working-Copy-Commit und macht
dann doch einen Snapshot.
Der Befehl nimmt nur Kennungen an, wie \type{op log} sie schreibt; Namen wie \type{main}
löst er nicht auf.

\subsection{Auf einen Pfad beschränken}

Ein Pfad am Ende des Aufrufs wählt eine Datei oder ein Verzeichnis samt allem darunter:

\starttyping
$ stemma diff b3:2037…df79 b3:9d8a…4c77 src
$ echo $?
0
$ stemma diff b3:2037…df79 b3:9d8a…4c77 nope.txt
stemma: `nope.txt` is not a file or directory on either side of the diff (…)
$ echo $?
1
\stoptyping

Unter \type{src} hat sich zwischen den Commits nichts geändert, der Diff ist leer, und der Befehl
endet mit 0. Der Pfad \type{nope.txt} kommt auf keiner Seite vor.
Das ist ein Fehler und kein leerer Diff, damit ein Tippfehler auffällt (Exit-Code 1).
Die Meldung ist hier gekürzt; sie nennt auch \type{track} als Abhilfe für eine Datei,
die nicht verfolgt wird.
Der Pfad beschränkt allein die Ausgabe. Mit dem Working-Copy-Commit als Seite speichert der
Snapshot weiter alle verfolgten Dateien, und das Log sieht keinen Unterschied zum Aufruf ohne Pfad.

\subsection{Den Diff in einem Skript prüfen}

Ein Skript will wissen, ob sich etwas geändert hat, und nicht den Text lesen.
Dafür gibt es \type{--exit-code}:

\starttyping
$ stemma diff --exit-code
$ echo $?
0
$ stemma diff --exit-code b3:2037…df79 b3:9d8a…4c77 > /dev/null
$ echo $?
4
\stoptyping

Der erste Aufruf vergleicht den leeren Working-Copy-Commit mit seinem Elternteil und endet mit 0.
Der zweite findet einen Unterschied und endet mit 4.
Die 4 ist kein Fehler: Git nimmt dafür 1, doch hier steht 1 für Fehler, die ein Skript von einem
Unterschied trennen muss (ADR-018). Mit Pfaden zählen nur die gewählten Dateien.
Die Option selbst schreibt nichts. Ein Snapshot entsteht nur, wenn eine Seite des Aufrufs der
Working-Copy-Commit ist.

\subsection{Die Oberfläche im Terminal}

Mit \type{--tui} blättert man durch die Dateien und Hunks, statt den Text zu lesen.
Die Oberfläche ist schreibgeschützt (ADR-016). Auf einem Terminal macht sie denselben Snapshot
wie der Text, danach schreibt sie nichts. Ohne Terminal auf der Standardausgabe endet der Befehl
mit Exit-Code 1, bevor er etwas liest oder speichert:

\starttyping
$ stemma diff --tui | cat
stemma: `diff --tui` needs a terminal on standard output; use `diff` without `--tui` for scripts
\stoptyping

Für Skripte und Agenten gilt deshalb der Text ohne \type{--tui}.

\subsection{Farbe und Pager gegenüber einer Pipe}

Auf einem Terminal färbt der Befehl die Ausgabe und gibt sie an einen Pager weiter.
Geht die Ausgabe in eine Pipe oder eine Datei, ist sie reiner Text ohne Farbe und ohne Pager.
Damit liest ein Skript denselben Text, den die Beispiele dieses Kapitels zeigen.
Farbe und Pager schreiben nichts ins Repository. Ein Pager, der früher endet, ist kein Fehler.
Welche Variablen den Pager wählen, steht in der Befehlsreferenz.

\subsection{Binärdateien}

Eine Datei mit einem NUL-Byte in den ersten 8000 Bytes gilt als binär.
Der Diff zeigt dann keinen Inhalt, sondern eine Zeile mit beiden Seiten.
Das Beispiel stammt aus einem eigenen kleinen Repository, nicht aus der Sitzung:
Es verfolgt und committet \type{blob.bin} und ändert danach ein Byte darin.
Seine IDs sind die des Laufs.

\starttyping
$ stemma diff
diff blob.bin
Binary files b3:c76d…402f blob.bin and working-copy b3:39f0…cdc7 blob.bin differ
\stoptyping

Mit \type{--exit-code} endet der Befehl auch bei Binärdateien mit 4.
Der Snapshot schreibt wie bei Text einen Eintrag \type{snapshot}, weil sich eine verfolgte Datei
geändert hat. Ist nur eine der beiden Fassungen binär, gilt die ganze Datei als binär.

\subsection{Was jeder Ablauf ändert}

\in{Tabelle}[tab:diff-effects] fasst zusammen, was die Abläufe in den vier Bereichen ändern.
Das Arbeitsverzeichnis ändert nie ein Diff. Ein Eintrag entsteht nur durch den Snapshot,
und der nur bei einer geänderten verfolgten Datei.

\startplacetable[title={Was ein Diff ändert},reference=tab:diff-effects]
\startxtable[frame=off]
\startxrow[topframe=on,bottomframe=on]
\startxcell {\bf Aufruf} \stopxcell
\startxcell {\bf Arbeitsverzeichnis} \stopxcell
\startxcell {\bf Working-Copy-Commit} \stopxcell
\startxcell {\bf Zeiger} \stopxcell
\startxcell {\bf Operation Log} \stopxcell
\stopxrow
\startxrow
\startxcell \type{diff}, Datei geändert \stopxcell
\startxcell unverändert \stopxcell
\startxcell neuer Snapshot \stopxcell
\startxcell Zeiger des Working-Copy-Commits bewegt \stopxcell
\startxcell ein Eintrag \type{snapshot} \stopxcell
\stopxrow
\startxrow
\startxcell \type{diff}, nichts geändert \stopxcell
\startxcell unverändert \stopxcell
\startxcell unverändert \stopxcell
\startxcell unverändert \stopxcell
\startxcell kein Eintrag \stopxcell
\stopxrow
\startxrow
\startxcell \type{diff} mit zwei Commits \stopxcell
\startxcell unverändert \stopxcell
\startxcell unverändert \stopxcell
\startxcell unverändert \stopxcell
\startxcell kein Eintrag \stopxcell
\stopxrow
\startxrow
\startxcell \type{diff} mit Pfad \stopxcell
\startxcell unverändert \stopxcell
\startxcell wie ohne Pfad \stopxcell
\startxcell wie ohne Pfad \stopxcell
\startxcell wie ohne Pfad \stopxcell
\stopxrow
\startxrow[bottomframe=on]
\startxcell \type{diff --tui} ohne Terminal \stopxcell
\startxcell unverändert \stopxcell
\startxcell unverändert \stopxcell
\startxcell unverändert \stopxcell
\startxcell kein Eintrag \stopxcell
\stopxrow
\stopxtable
\stopplacetable

Nach einem Absturz im Snapshot gilt, was \in{Kapitel}[chap:phase1] für den Snapshot sagt.
Der Diff schreibt sonst nichts, ein erneuter Aufruf ist deshalb jederzeit sicher.

\section{Offene Fragen}

Die folgenden Fragen sind offen.

\startitemize
\item Eine Option oder Variable für die Farbe, etwa \type{NO_COLOR}, und die Wahl des Pagers
über die Einstellungen von Git.
\item Eine JSON-Form für \type{diff} und die übrigen Befehle ohne \type{--json},
mit der Schnittstelle der Phase 6 (ADR-018).
\stopitemize
